%% cover_letter_emse.tex
%% Compile separately: pdflatex cover_letter_emse.tex

\documentclass[9pt]{letter}
\usepackage[utf8]{inputenc}
\usepackage[T1]{fontenc}
\usepackage[margin=0.65in]{geometry}

\begin{document}

\begin{letter}{Editor-in-Chief\\
Empirical Software Engineering}

\opening{Dear Editor,}

\textbf{Subject: Submission of manuscript}\\
\textbf{``Recovery-Channel Prompt Injection in Self-Healing LLM Agent Frameworks''}

We are pleased to submit the above-referenced manuscript for consideration by
\emph{Empirical Software Engineering}, specifically the special issue
\emph{Agentic Software Engineering: The Rise of AI Teammates}.

\textbf{Summary.}
This paper studies recovery-channel prompt injection: externally originated
failure content is captured by a self-healing agent and re-injected into a
subsequent LLM turn under a recovery frame that can be interpreted as
authoritative.  We combine systematic multiple-case source-code analysis with
native-framework experiments.  In the canonical Reflexion matrix, five model
aliases produce attack success rates from 0.174 to 0.796 while all matched
benign controls are zero.  Native AutoGen and Aider executions provide
additional framework evidence, and supplementary Reflexion contrast, defense-aware,
and legitimate-recovery screens delimit what can and cannot be inferred at the
completed sample sizes.  A native Reflexion TrustOrigin port reduces attack
success from 0.776 to 0.548 with zero matched benign false positives.  The
paper's design implication is to pair provenance
handling with structural authorization checks at recovery sinks.

The manuscript directly fits three themes of the special issue.  It studies
\emph{Failure Patterns and Risks} through a cross-framework taxonomy of
recovery-path trust gaps and security vulnerabilities; \emph{Architectures for
Agents} through the source-level analysis of retry, tool-error, and recovery
transport paths; and \emph{Testing and Evaluation} through native-framework
execution, matched benign controls, and framework/model/payload-stratified
measurements.

We do not claim a new underlying language-model failure mode.  The contribution
is an empirical software-engineering characterization of a distinct framework
transport path and trust boundary: externally influenced failure content is
captured, reframed, and re-injected by recovery infrastructure.  The unit of
analysis is the fixed-revision agent recovery implementation, and the results
are reported for maintainers and designers of agent frameworks rather than as a
new model-level attack technique.

The study is an empirical software-engineering contribution: systematic
multiple-case source analysis with qualitative coding, fixed framework/model/
payload strata, matched controls, raw logs, and explicit uncertainty reporting.
Provider aliases and endpoint metadata are disclosed, while the lack of
immutable hosted checkpoints is stated as a reproducibility limitation.

\textbf{Fit with the special issue and EMSE.}
The replication package preserves canonical manifests, raw logs, endpoint
verification, model provenance, and the numerical audit used for the tables.
A fixed local Qwen2-family \texttt{qwen:7b} run is included as an exploratory
open-model sensitivity anchor (100 valid trials). The native Reflexion
TrustOrigin contrast is the primary mitigation result; defense-aware and
legitimate-recovery screens are supplementary, while structural
reauthorization is a design implication rather than a separate empirical RQ.

We confirm that this manuscript has not been published and is not under
consideration elsewhere.  We look forward to the editors' and reviewers'
feedback.

\closing{Sincerely,\\[1em]
Mingyuan Xie (jhj208436@gmail.com)\\
Zhengxun Wu (2276815088@qq.com)\\
School of Information Management, Qingdao University of Technology\\
Qingdao, Shandong, China}

\end{letter}

\end{document}
